\section{An exact ordinary full-chain fiber from a two-dimensional branch}
\label{sec:cubic-exact-ordinary-fiber}

The Frobenius convention in this section is
\(\epsilon(\operatorname{Frob}_l)=l^{-1}\) for geometric Frobenius.
The normalized two-dimensional parameter has determinant \(\eta_K\);
its unitary local Frobenius roots have product \(\eta_K(l)\).
Thorne defines
\(\operatorname{rec}^T(\pi)=\operatorname{rec}
(\pi|\det|^{(1-n)/2})\), printed p.~6.
Thus the rank-three determinant attached to a unitary parameter
\((\chi^j V)\oplus1\) is
\[
       \epsilon^{-1}\bigl((\chi^j V)|_{G_F}\oplus1\bigr).
       \tag{G78.1}
\]
The argument uses the literal definition and the Hecke polynomial
\(P_w\), not the common Tate twist in the display following the
proof of Thorne's Theorem~3.3.

\begin{lemma}[AC-G78-module: the actual universal ordinary module]
\label{lem:G78-universal}
Fix a sufficiently small tame level \(K^p\) in Thorne's Section~6,
and put \(\Lambda_G=\mathcal O[[\mathfrak t_p]]\).
There is a finite free \(\Lambda_G\)-module \(\mathscr M(K^p)\),
with its actual commuting away Hecke action, such that for every
continuous character \(\lambda:\mathfrak t_p\to\mathcal O'^\times\)
and finite coefficient extension \(\mathcal O'/\mathcal O\),
\[
 \mathscr M(K^p)\otimes_{\Lambda_G,w_0\lambda^{-1}}\mathcal O'
       \simeq H_\lambda(K^p;\mathcal O')
       \tag{G78.2}
\]
equivariantly for the away Hecke operators.
No residual irreducibility hypothesis is used.
\end{lemma}

\begin{proof}
Write \(D_d=\mathfrak t_p/\mathfrak t_p(d)\) and
\(M_d=M(K^p(d,d),\mathcal O)^{\rm ord}\), using the weight-zero
spaces on printed pp.~33--35 of Thorne.
The full space at level \(d\) is free over \(\mathcal O[D_d]\):
this is the sufficiently-small-level statement preceding (6.1).
Ordinary projection is a commuting integral direct-summand
projection, by Lemma~6.1.  Hence \(M_d\) is projective over
\(\mathcal O[D_d]\), and is free because this ring is local
(\(D_d\) is a finite abelian \(p\)-group).
The trace, together with the ordinary isomorphism when only the
lower-triangular depth is increased, gives
\[
 M_{d+1}\otimes_{\mathcal O[D_{d+1}]}\mathcal O[D_d]
                      \simeq M_d.                    \tag{G78.3}
\]
This is precisely (6.1), applied first at common lower depth
\(d+1\), followed by the level-depth ordinary isomorphism cited
immediately after that equation.  In particular the transition
is surjective.  All ranks over the respective group rings are
the same finite integer \(m\).
Choose a basis at \(d=1\), and lift it successively through
(G78.3).  Nakayama's lemma over the local group rings shows that
each lifted tuple is a basis.  Taking inverse limits along these
trace maps gives
\(\mathscr M=\varprojlim_d M_d\simeq\Lambda_G^m\).
Every away Hecke operator commutes with these traces.
If \(I_{d,d+1}\) denotes the natural inclusion, with the ordinary
lower-depth identification understood, the compatibility is literally
\[
 I_{d,d+1}(\operatorname{Tr}_{d+1,d}m)
       =\sum_{a\in\ker(D_{d+1}\to D_d)}a\,m.
\]
Thus the inverse transition and the direct inclusion differ by
the stated kernel norm, with no scalar normalization.

For completeness, specialization uses a norm, not division by
\(|D_d|\).  If a character \(\xi:D_d\to(\mathcal O'/\varpi'^c)^\times\)
is given, then on a free group-ring module the twisted norm
\(\sum_{a\in D_d}\xi(a)^{-1}a\) induces an isomorphism from
coinvariants modulo the relations \(a=\xi(a)\) to the
\(\xi\)-eigenspace.  This can be checked on the basis of a regular
group ring: both modules are free of rank one and the image of
\(1\) is a generator of the eigenspace.
For a continuous \(\lambda\), choose \(d\) so that
\(w_0\lambda^{-1}\bmod\varpi'^c\) factors through \(D_d\).
The preceding norm isomorphism, the trace compatibility
(G78.3), and the identification of
\(H_\lambda/\varpi'^c\) in the proof of Proposition~6.2 give
(G78.2) modulo \(\varpi'^c\).
The norm maps are compatible when \(d\) is increased:
the norm over a larger group is the product of the norm over
the kernel and the lifted norm over the quotient.
Taking the inverse limit in \(c\) proves (G78.2).
These are the same finite free group-ring and twisted-norm
calculations used in the proof of Proposition~6.7, printed
pp.~39--40.
\end{proof}

\begin{theorem}[AC-G78: an actual full-chain ordinary fiber]
\label{thm:G78}
Let \(p\ge5\), \(K=\mathbf Q(\sqrt{-3})\), and let
\(\chi\) be the fixed anticyclotomic cubic character of conductor
\(269\).  Let \(f_0\) be an ordinary cuspidal eigenform of weight
\(2e+1\), \(e\ge1\), nebentype \(\eta_K\), with underlying newform level prime
to \(p\), and choose its ordinary \(p\)-stabilization.  Assume its restriction to \(G_K\) is irreducible.
A fixed quadratic character \(\tau\), trivial or of conductor
\(p\), may be included: the starting normalized endpoint
parameter is the unitary normalization of \(f_0\), tensored
with \(\tau\).
Its determinant is \(\eta_K\).
Assume the inherited residual shape
\[
 \overline V^{\rm ss}|_{G_K}=
       \overline\gamma\oplus\overline\gamma^{-1},
 \qquad \gamma=\psi\,\epsilon^e,
\]
where \(\psi\) is a finite rational character and its inertia at \(p\)
is the quadratic character \(\tau\); the corresponding global
two-dimensional residual characters are
\(\eta_K\overline\gamma,\overline\gamma^{-1}\).
This is the actual old Type II residual shape.
The theorem also applies after the already proved good-twist
construction has supplied this \(f_0,\tau\).

Take the real quadratic host of AC-G77:
\(F=KK'\), \(F^+=\mathbf Q(\sqrt{3r})\),
\(K'=\mathbf Q(\sqrt{-r})\), with the fresh prime \(r\) chosen
there, and use \(F_0=K'\) in Thorne.
Include every old ramified prime and \(p\) in its prescribed
splitting set.
Fix \(j\in\{1,2\}\).

Here is the exact sign premise.
On inertia at \(p\), the normalized ordinary branch has
semisimple characters
\(\tau\,\omega^e\langle\epsilon\rangle^{w}\) and
\(\tau\,\omega^{-e}\langle\epsilon\rangle^{-w}\),
where \(w=e\) is the starting weight.
Let
\[
                       s=\tau(-1)(-1)^e\in\{1,-1\}.
\]
If \(s=-1\), no extra scalar premise is required.
If \(s=1\), assume there is
\[
 h_0\in G_F,\qquad
 \overline V(h_0)=-I_2,\quad
 \overline\epsilon(h_0)=1,\quad \chi(h_0)=1.
 \tag{G78.4}
\]
The residual equality is a semisimple equality, or equivalently
trace \(-2\) and determinant \(1\); the proof below only uses the
pseudotrace.
It is the actual restriction in (G78.4), not an assertion about
the order of a character before restriction to the host.

Then there exist a rational prime \(q\), split completely in
\(F\), with
\[
 q\equiv-1\pmod p,\qquad \chi(\operatorname{Frob}_q)=1,
 \tag{G78.5}
\]
a finite coefficient extension \(E'/E\), a continuous ordinary
weight \(\lambda_q\), and an actual characteristic-zero Hecke
point on \(H_{\lambda_q}(K^p;E')\), at a common tame level
hyperspecial at all places above \(q\), whose source-normalized
Frobenius polynomial at the chosen \(q\)-place is
\[
                       (X-1)(X-q)(X-q^2).             \tag{G78.6}
\]
The residual system is the one attached to the original
\((\chi^j V)\oplus1\), with the common Tate normalization (G78.1).
In particular AC-G75 gives an actual ordinary Steinberg
constituent at the original symmetric classical weight and
the fixed away-\(v_0\) tame level.
The source genericity hypothesis can be fulfilled for this
cubic residual system.
In particular, for the inherited actual \(269\)-coefficient
Type II endpoints, the conclusion holds for every odd \(e\):
AC-G68 supplies the old modular endpoint, AC-K29-nonCM supplies
irreducibility, and the good-twist construction supplies \(f_0,\tau\).
Since \(269\equiv1\pmod4\), both allowed \(\tau\) have
\(\tau(-1)=1\), so the sign is \(-1\).
Only the good-twist proof is used here; no Euler-system conclusion
at \(269\) is imported.

No bound on \(q\), no original-level descent, and no
per-curve containment in the modified cubic deformation ring
is asserted.
\end{theorem}

\begin{proof}
We give the branch construction, the exact root argument, and
the actual Hecke-support passage separately.

\emph{The ordinary branch and its normalization.}
The parabolic ordinary Hecke algebra is finite projective over
its one-variable weight algebra.  This is KLZ,
\emph{Rankin--Eisenstein classes and explicit reciprocity laws},
arXiv:1503.02888v5, Proposition~7.2.1 and the paragraph preceding
Theorem~7.2.3, printed pp.~35--37.
Choose a minimal prime below the point of \(f_0\), and let \(R\)
be the resulting complete local domain.
It is finite and torsionfree over the chosen weight component
\(\mathcal O[[W]]\); no claim that the entire old-level Hecke
algebra is reduced is needed.
Use the half-weight coordinate for which the universal
character at \(u=1+p\) is \(1+W\); at classical weight \(2w+1\)
it is \(u^w\).
Multiplication of the usual weight coordinate by \(2\) is an
automorphism because \(p\ne2\).

The branch has a continuous rank-two determinant and trace.
This is the ordinary quotient of the standard Hecke-valued
determinant; an explicit source with no residual irreducibility
assumption is Pan, \emph{On locally analytic vectors of the
completed cohomology of modular curves}, Section~6.1.1, published printed p.~67, which explicitly constructs
its continuous determinant with polynomial
\(X^2-l^{-1}T_lX+l^{-1}S_l\).
Equivalently it follows by the classical determinant gluing
used in Thorne's Proposition~6.4(1).
More explicitly, if \(D_{\rm Pan}\) denotes this determinant,
first twist by \(\epsilon^{-1}\) to obtain the usual classical
Frobenius trace \(a_l\), then by the analytic character
\(\omega^{e}\langle\epsilon\rangle^{w}\), and by \(\tau\).
The first two twists combine to
\(\omega^{e-1}\langle\epsilon\rangle^{w-1}\).
On the dense progression \(w\equiv e\pmod{p-1}\) this is
exactly \(\epsilon^{w-1}\).
Thus the normalized determinant is
\(\tau\,\omega^{e-1}\langle\epsilon\rangle^{w-1}D_{\rm Pan}\),
where multiplying a determinant means twisting the underlying
rank-two polynomial law.
Its classical good-prime trace is \(\tau(l)a_l/l^w\).
The resulting
pseudotrace \(T:G_{\mathbf Q}\to R\) has determinant
\(\eta_K\).
At a classical point of weight \(2w+1\), its good-prime
unitary trace is \(a_l(f)/l^w\), multiplied by \(\tau(l)\).
This specifies the normalization independently of whether
the unnormalized cohomology realizes the dual Deligne
representation.
On \(G_F\) the determinant is \(1\).

The classical points of weights \(2w+1\), with \(w>0\) integral
and \(w\equiv e\pmod{p-1}\), are Zariski dense on this branch.
The ordinary control assertion here is KLZ Section~7.3.
At these weights \(2w+1\ge3\), with fixed finite nebentype
\(\eta_K\), the ordinary specialization is the ordinary
stabilization of a form unramified at \(p\): a Steinberg
newform at \(p\) has \(U_p\)-slope \((2w-1)/2>0\), and
a ramified principal series with trivial inertial central
character cannot have exactly one unramified inducing
character.  Thus neither alternative supplies a unit
\(U_p\)-eigenvalue.
The non-Eisenstein hypotheses of Theorem~7.2.8 are not used.
The classically specialized two-dimensional trace is ordinary,
so on \(p\)-inertia the trace has the two characters stated in
the theorem.  Their trace identity holds in \(R\), since it
holds at the dense classical points.

The restriction of a dense set of these classical forms to
\(G_K\) remains irreducible.  Indeed, irreducibility of the
starting restriction means that the starting representation
has no quadratic self-twist by \(\eta_K\).
There is a group element \(a\) for which the analytic function
\(T(a)-\eta_K(a)T(a)\) is nonzero at that point.
Its nonvanishing is a Zariski-open condition, and excludes
that quadratic self-twist.
Quadratic base change of a cuspidal \({\rm GL}_2\) form fails
to be cuspidal only in this self-twist case.
The further extension from \(K\) to \(F\) cannot introduce
reducibility: all forms have level prime to the fresh \(r\),
whereas the new quadratic character is ramified there.
An unramified local parameter cannot equal its ramified
scalar quadratic twist.
Thus a dense set of these classical points gives the
cuspidal endpoint base change required in Thorne's
Theorem~4.5.

\emph{A scalar witness needs no free Galois lattice.}
Suppose \(s=1\) and take \(h_0\) as in (G78.4).
Put \(t_0=T(h_0)\) and
\[
 A=R[X]/(X^2-t_0X+1).
\]
This is finite free over \(R\).  In \(A/\mathfrak m_RA\),
the element \(-X\) is unipotent with square-zero deviation
from \(1\).
Hence \((-X)^p\equiv1\pmod{\mathfrak m_RA}\).
Repeated \(p\)-th powering sends \(1+\mathfrak m_RA\) to \(1\)
in the \(\mathfrak m_R\)-adic topology, and therefore
\(X^{p^n}\to-1\).
The two-dimensional Cayley--Hamilton trace identity
\[
 T(h_0^2a)-T(h_0)T(h_0a)+T(a)=0
 \tag{G78.7}
\]
holds for every \(a\); it may be checked at all classical
points and then in the domain \(R\).
Applying the corresponding polynomial recurrence shows
\[
                     T(h_0^{p^n}a)\longrightarrow -T(a)
\]
\(\mathfrak m_R\)-adically, for every \(a\).
Choose a convergent subsequence of \(h_0^{p^n}\) in the
compact Galois group.  Its limit \(h\in G_F\) satisfies
\[
 T(ha)=-T(a)\quad\hbox{for all }a,\qquad
 \epsilon(h)=1,\qquad\chi(h)=1.
 \tag{G78.8}
\]
Here the cyclotomic assertion uses
\(\epsilon(h_0)\in1+p\mathbf Z_p\), and the cubic assertion
uses \(\chi(h_0)=1\).
Thus (G78.8) is a scalar identity for the actual pseudotrace;
no free rank-two representation over \(R\), residual
irreducibility, or big-image theorem has been assumed.

Choose \(t\in I_p\cap G_F\) with
\(\epsilon(t)=-u\) and \(\chi(t)=1\), \(u=1+p\).
At \(p\ne269\) the cubic character is unramified.
At \(p=269\), its tame order \(3\) and the cyclotomic tame
order \(268\) are coprime, so the two prescribed values can
still be chosen simultaneously; the wild cyclotomic value
is unrestricted.
The host is unramified at \(p\), locally equal to \(K_p\)
when it is not split, so it does not remove this inertia.
The ordinary trace formula gives
\[
 T(t)=s\bigl((1+W)+(1+W)^{-1}\bigr).
\]
If \(s=-1\), put \(g=t\).
If \(s=1\), put \(g=ht\), with \(h\) from (G78.8).
In both cases,
\[
 T(g)=-\bigl((1+W)+(1+W)^{-1}\bigr),\quad
 \epsilon(g)=-u,\quad \chi(g)=1.                       \tag{G78.9}
\]

\emph{An exact nearby Frobenius equation.}
There is a point of the generic fiber of \(R\) above \(W=p\),
since \(R\) is finite integral and dominant over
\(\mathcal O[[W]]\).
Normalize that analytic curve and choose a point \(x\)
above \(W=p\), over a finite coefficient extension.
Its local analytic parameter can be written \(z\), with
\[
                        W-p=z^d b(z),\qquad b(0)\ne0.
\]
Consider
\[
 F_g(W)=T(g)-\bigl(\epsilon(g)+\epsilon(g)^{-1}\bigr)
 =-\bigl((1+W)+(1+W)^{-1}\bigr)+u+u^{-1}.
\]
It vanishes at \(W=p\), with nonzero derivative
\(-1+u^{-2}\).
Consequently \(F_g\) has a zero of order exactly \(d\)
in the parameter \(z\).
Choose a sufficiently small closed parameter disc on which
its leading term strictly dominates the higher terms
on the boundary.

The map from the complete local ring \(R\) to the bounded
analytic functions on this disc is continuous for the
\(\mathfrak m_R\)-adic topology.
To see this, integrality over \(\mathcal O[[W]]\) first bounds
the images of the elements of \(R\) on the chosen integral disc:
a root of a monic polynomial with coefficients of norm at most
one also has norm at most one.
Moreover \(R\) is finite over
\(\mathcal O[[W]]\), and a power of its maximal ideal is
contained in \((\varpi,W)R\).  On the chosen disc both
\(\varpi\) and \(W\) have norm less than one.
Hence a sufficiently high power of \(\mathfrak m_R\)
has arbitrarily small uniform norm.
Continuity of \(T\) and \(\epsilon\) therefore allows a
finite Galois quotient such that any Frobenius conjugate
sufficiently close to \(g\) gives a function
\[
 F_q=T(\operatorname{Frob}_q)-(q+q^{-1})
\]
uniformly closer to \(F_g\) than its boundary norm.
Choose \(q\) by Chebotarev, avoiding all fixed bad primes.
The quotient also retains \(G_F\), \(\chi=1\), and
\(\overline\epsilon=-1\), so \(q\) splits completely in \(F\)
and satisfies (G78.5).

The elementary nonarchimedean Weierstrass theorem now gives
a zero \(x_q\) of \(F_q\) inside the disc, after a finite
coefficient extension.
Explicitly, rescale the disc to the unit disc and divide
by the leading boundary coefficient.  Its distinguished
polynomial has degree \(d\); a perturbation of strictly
smaller Gauss norm has the same distinguished degree,
and therefore has \(d\) roots in the disc, counted with
multiplicity.
This does not assume that the weight map is etale.
At \(x_q\), determinant \(1\) and trace \(q+q^{-1}\) imply
that the normalized two-dimensional Frobenius roots
are exactly \(q,q^{-1}\).
Since \(\chi(\operatorname{Frob}_q)=1\), the rank-three
parameter before the common Tate factor has roots
\(q,1,q^{-1}\).
Formula (G78.1) gives (G78.6).

\emph{The root lies in the actual ordinary Hecke support.}
This step is algebraic; it does not approximate an arbitrary
algebraic weight by integer weights.
For each dense classical point considered above, put
\[
 \pi_w=(\operatorname{BC}_F(f_w)^{\rm unit}\otimes
                    \tau\chi^j)\boxplus1 .
\]
It is conjugate self-dual, regular algebraic, and ordinary.
Its unitary infinity exponents are \(w,0,-w\), so its
algebraic weight is
\(\mu(w)=(w-1,0,1-w)\) at each embedding.
The even multiplicity condition of Thorne's
Theorem~4.5 holds by AC-G77.
Every ramified place splits over \(F^+\).
The endpoint is cuspidal on the dense open subset already
specified.  Thus Theorem~4.5 and Proposition~6.5 give the
actual ordinary Hecke system on the definite group.

A common tame level \(K^p\) can be chosen for all these
points: the tame conductors of the two-dimensional family
and of \(\chi,\tau\) are bounded, and fixed principal
congruence subgroups are contained in their newvector
subgroups.  Choose it sufficiently small, and hyperspecial
at every place over the good prime \(q\).
At \(p\), the fixed twists are tame, hence trivial on
\(\mathfrak t_p\).  Therefore the character
\(\lambda(w)\) in Proposition~6.5 is just the restriction
of \(\mu(w)\).
In the universal branch it is the continuous \(R^\times\)-
valued character
\[
 \lambda_R(t)=
 \frac{\nu_R\bigl(\prod_{v\mid p}
             N_{F_{\widetilde v}/\mathbf Q_p}(t_{v,1}/t_{v,3})\bigr)}
      {\prod_{v\mid p}
             N_{F_{\widetilde v}/\mathbf Q_p}(t_{v,1}/t_{v,3})},
 \tag{G78.10}
\]
where \(\nu_R(u)=1+W\) on \(1+p\mathbf Z_p\).
The local fields at \(p\) are unramified, so this formula
is the product of the parallel embedding characters.
At every classical point it gives exactly the source
ordinary weight.

Base change the finite free module in Lemma
\ref{lem:G78-universal} along
\(w_0\lambda_R^{-1}:\Lambda_G\to R\).
Let \(B\) be the \(R\)-algebra generated by its actual
away Hecke operators inside its endomorphism algebra.
It is a finite \(R\)-algebra.
Define the desired values of the unramified Hecke
generators by the characteristic polynomial of
\(\epsilon^{-1}((\chi^j T)\oplus1)\).
These values lie in \(R\), including the inverse of
the final coefficient.
Every relation among these actual operators evaluates
to zero at each dense classical point, because of
Theorem~4.5 and Proposition~6.5.
It is therefore zero in the domain \(R\).
Thus this assignment defines an \(R\)-algebra map
\[
                               B\longrightarrow R.    \tag{G78.11}
\]

Specialize (G78.11) at \(x_q\).
The resulting character factors through the actual
Hecke image on the specialized finite free module.
Indeed, if an element \(b\in B\) acts as zero after
specialization, its characteristic polynomial as an
endomorphism specializes to \(X^m\).
Cayley--Hamilton gives \(b^m\in\mathfrak m_{x_q}B\)
after extending to the coefficient field of \(x_q\).
Every specialized character therefore kills \(b\).
The Cayley--Hamilton assertion here is the polynomial-law
identity for the universal matrix representation
\(B\subseteq\operatorname{End}_R(\mathscr M_R)\).
It survives arbitrary residue-field and finite coefficient
base change, so it also applies to every element of the
base-changed algebra, not only to elements coming from \(B\).
By (G78.2), this specialized module is the actual
\(H_{\lambda_q}(K^p)\).
This proves the asserted Hecke occurrence and (G78.6).

Finally, the required generic residual prime exists.
Choose a Galois element fixing \(F\) with cyclotomic
residue \(-1\), rational endpoint residue \(\gamma=\pm1\),
and nontrivial cubic value.
One may start with \(p\)-inertia for the first two
conditions and then vary the cubic value in the kernel
of all rational abelian characters.
The latter variation is possible because the cubic
extension of \(K\) is not contained in an abelian
extension of \(\mathbf Q\); its intersection with the
abelian host and the rational character fields is \(K\).
The three normalized residual eigenvalues are then
\(\pm\zeta_3,1,\pm\zeta_3\).
Their ordered ratios are \(1\) or \(\pm\zeta_3^{\pm1}\),
none equal to \(-1\), since \(p\ne2,3\).
Chebotarev gives a good split prime with this property,
different from all previously chosen primes.
This is the literal genericity hypothesis of
Thorne's Theorem~5.2; distinct residual eigenvalues
are not required there.

The actual fiber (G78.6) now satisfies AC-G75.
Its conclusion supplies the asserted classical
ordinary Steinberg constituent at the original
weight.  All qualifications in the theorem have
been maintained throughout.
\end{proof}

\begin{remark}[Uncovered branches and later interfaces]
If the original rational finite character is principal
and \(e\) is even, the inertia sign is \(+1\) and
\(\overline V=\overline\epsilon^e\oplus
\overline\epsilon^{-e}\) on \(G_K\).
No element of \(\ker\overline\epsilon\) can have
residual value \(-I\).
Thus neither sign alternative above covers that branch.
A finite character killed by the chosen host cannot
supply (G78.4).
The theorem is a genuine occurrence statement on
its stated branch, not a universal ordinary level-
raising theorem for every residual datum.
Chebotarev and Weierstrass give no conductor-uniform
bound for the new \(q\).
The finite-type transfer of AC-G72 may be applied
only when its own projectivity and local hypotheses
are fulfilled, and neither that transfer nor the
present theorem proves the per-curve Hecke-ideal
containment needed for the original modified ring.
\end{remark}
